% problem-000327 10 芳香化合物同分异构
\section*{10. 芳香化合物同分异构}
$
\begin{array}{r l} & {\mathbf {A} \xrightarrow {\mathrm{NaOH} \text {溶液}} \xrightarrow {\text {酸化}} \mathbf {E} (\mathrm{C} _ {7} \mathrm{H} _ {6} \mathrm{O} _ {2}) + \mathbf {F}} \\ & {\mathbf {B} \xrightarrow {\mathrm{NaOH} \text {溶液}} \xrightarrow {\text {酸化}} \mathbf {G} (\mathrm{C} _ {7} \mathrm{H} _ {8} \mathrm{O}) + \mathbf {H}} \\ & {\mathbf {C} \xrightarrow {\mathrm{NaOH} \text {溶液}} \xrightarrow {\text {酸化}} \mathbf {I} (\text {芳香化合物}) + \mathbf {J}} \\ & {\mathbf {D} \xrightarrow {\mathrm{NaOH} \text {溶液}} \mathbf {K} + \mathrm{H} _ {2} \mathrm{O}} \end{array}
$

\subsection*{10-1}
写出A、B、C和D的分⼦式。

\subsection*{10-2}
画出A、B、C和D的结构简 $\updownarrow \updownarrow _ { \updownarrow }$

\subsection*{10-3}
A和D分别与NaOH溶液发⽣了哪类反应？

\subsection*{10-4}
写出H分⼦中官能团的名称。

\subsection*{10-5}
现有如下溶液：HCl、 $\mathrm { H N O _ { 3 } \setminus \ N H _ { 3 } { \cdot } H _ { 2 } O }$ 、NaOH、 $\mathrm { N a H C O _ { 3 } } ,$ 饱和 $\mathrm { B r } _ { 2 }$ ⽔、 $\mathrm { F e C l _ { 3 } }$ 和 $\mathrm { N H _ { 4 } C l _ { o } }$ 从中选择合适试剂，设计⼀种实验⽅案，鉴别E、G和 $\mathbf { I } _ { \circ }$
